\documentclass[12pt]{article}
\usepackage{amsmath, amssymb}
\usepackage[a4paper, margin=1in]{geometry}
\usepackage{enumitem}

\title{\textbf{CSE400 Lecture 4 Scribe}\\
Joint Probability and Conditional Probability}
\author{Course: CSE400 -- Fundamentals of Probability in Computing\\
Instructor: Dhaval Patel, PhD}
\date{January 15, 2026}

\begin{document}

\maketitle

\section{Lecture Outline}

The lecture develops probability theory from basic foundations and then introduces
joint probability.

Topics covered:

\begin{itemize}
    \item Introduction to Probability Theory
    \item Experiments, Sample Space, and Events
    \item Axioms of Probability
    \item Corollaries and Propositions from Probability Axioms
    \item How to Assign Probability
    \begin{itemize}
        \item Classical Approach
        \item Relative Frequency Approach
    \end{itemize}
    \item Joint Probability
    \begin{itemize}
        \item Motivation, Notation, and Concepts
        \item Example 1: Card Deck Example
        \item Example 2: Costume Party Example
    \end{itemize}
    \item Conditional Probability
    \begin{itemize}
        \item Motivation, Notation, and Concepts
        \item Example 3: Cards Without Replacement
        \item Example 4: Game of Poker
        \item Example 5: The Missing Key
    \end{itemize}
\end{itemize}

\section{Introduction to Probability Theory}

Probability theory begins with defining experiments, outcomes, events, and sample
spaces before assigning probabilities.

\section{Experiments, Sample Space, and Events}

\subsection{Experiment}

\textbf{Definition:}  
An Experiment ($E$) is a procedure we perform that produces some result.

\textbf{Example:}  
Tossing a coin five times is denoted as:

\[
E_5
\]

\subsection{Outcome}

\textbf{Definition:}  
An Outcome ($\xi$) is a possible result of an experiment.

\textbf{Example:}  
One outcome of $E_5$ is:

\[
\xi_1 = HHTHT
\]

\subsection{Event}

\textbf{Definition:}  
An Event is a certain set of outcomes of an experiment.

\textbf{Example:}  
Event $C$ for experiment $E_5$:

\[
C = \{\text{all outcomes consisting of an even number of heads}\}
\]

\subsection{Sample Space}

\textbf{Definition:}  
The Sample Space ($S$) is the collection or set of all possible distinct outcomes of an
experiment.

Key properties:

\begin{itemize}
    \item \textbf{Mutually Exclusive:} outcomes cannot occur together
    \item \textbf{Collectively Exhaustive:} outcomes cover all possibilities
\end{itemize}

\textbf{Example: Flip a coin}

\begin{itemize}
    \item Heads or tails, not both
    \item Nothing else is possible
\end{itemize}

Sample space may be:

\begin{itemize}
    \item Discrete
    \item Countably infinite
    \item Continuous
\end{itemize}

Examples listed:

\begin{itemize}
    \item Flipping a fair coin once
    \item Rolling a cubical die
    \item Rolling two dice
    \item Flip until a tail occurs
    \item Random number generator in $[0,1)$
\end{itemize}

\section{Axioms of Probability}

After defining events, the next step is assigning probabilities.

\subsection{Probability Definition}

Probability is:

\begin{itemize}
    \item A measure of the likelihood of various events
    \item OR a function of an event that produces a numerical quantity measuring likelihood
\end{itemize}

\subsection{Axioms}

Axioms are self-evident and require no proof.

\subsubsection*{Axiom 1 (Bounds)}

For any event $A$:

\[
0 \le Pr(A) \le 1
\]

\subsubsection*{Axiom 2 (Certain Event)}

If $S$ is the sample space:

\[
Pr(S) = 1
\]

\subsubsection*{Axiom 3 (Additivity)}

If $A \cap B = \emptyset$:

\[
Pr(A \cup B) = Pr(A) + Pr(B)
\]

For infinitely many mutually exclusive sets $A_i$:

\[
Pr\left(\bigcup_{i=1}^{\infty} A_i\right) = \sum_{i=1}^{\infty} Pr(A_i)
\]

\section{Corollary from Probability Axioms}

\textbf{Corollary 2.1}  
For a finite number $M$ of mutually exclusive sets:

\[
Pr\left(\bigcup_{i=1}^{M} A_i\right) = \sum_{i=1}^{M} Pr(A_i)
\]

Important slide notes:

\begin{itemize}
    \item Axiom 3 is equivalent to Corollary 2.1 when the sample space is finite
    \item The generality of Axiom 3 is necessary for infinite sample spaces
\end{itemize}

\section{Propositions from Probability Axioms}

\subsection*{Proposition 2.1 (Complement Rule)}

\[
Pr(A^c) = 1 - Pr(A)
\]

\subsection*{Proposition 2.2 (Subset Rule)}

If $A \subset B$:

\[
Pr(A) \le Pr(B)
\]

\subsection*{Proposition 2.3 (Union Rule)}

For any sets $A$ and $B$:

\[
Pr(A \cup B) = Pr(A) + Pr(B) - Pr(A \cap B)
\]

\subsection*{Proposition 2.4 (Inclusion--Exclusion Principle)}

For events $A_1, \dots, A_M$:

\[
Pr\left(\bigcup_{i=1}^{M} A_i\right)
= \sum_i Pr(A_i) - \sum_{i<j} Pr(A_iA_j) + \cdots + (-1)^{M+1}Pr(A_1A_2\cdots A_M)
\]

\section{Assigning Probabilities}

\subsection{Classical Approach}

Examples listed:

\begin{itemize}
    \item Coin flipping: compute $Pr(H)$, $Pr(T)$
    \item Dice rolling: probability of even number
    \item Pair of dice: probability sum equals five
\end{itemize}

\subsection{Relative Frequency Approach}

Example: Pair of dice, sum equals five.

Inference from slides:

\begin{itemize}
    \item Exact probability requires infinite repetitions
    \item Many random phenomena are not repeatable
\end{itemize}

\section{Joint Probability}

\subsection{Motivation}

All events are not mutually exclusive.

We are interested in the probability of intersection:

\[
A \cap B
\]

\subsection{Notation}

Joint probability is denoted as:

\[
Pr(A,B) \quad \text{or} \quad Pr(A \cap B)
\]

\subsection{Multiple Events}

For events $A_1, A_2, \dots, A_M$:

\[
Pr(A_1, A_2, \dots, A_M)
\]

\subsection{Calculation Approaches}

\subsubsection*{Approach 1: Classical}

Step-by-step reasoning:

\begin{enumerate}
    \item Express both events $A$ and $B$ using atomic outcomes
    \item Write the set of atomic outcomes common to both
    \item Calculate probabilities of those outcomes
\end{enumerate}

\subsubsection*{Approach 2: Relative Frequency}

Let $n_{AB}$ be the number of times $A$ and $B$ occur simultaneously in $n$ trials:

\[
Pr(A,B) = \lim_{n\to\infty} \frac{n_{AB}}{n}
\]

\section{Worked Examples}

\subsection{Example 1: Card Deck Example}

A deck has 52 cards.

Define:

\[
A = \{\text{Red card selected}\}
\]
\[
B = \{\text{Number card selected (Ace is a number card)}\}
\]
\[
C = \{\text{Heart card selected}\}
\]

Find:

\[
Pr(A), Pr(B), Pr(C), Pr(A,B), Pr(A,C), Pr(B,C)
\]

\subsubsection*{Solution}

Step 1: Individual Probabilities

\[
Pr(A) = \frac{26}{52}
\]

\[
Pr(B) = \frac{40}{52}
\]

\[
Pr(C) = \frac{13}{52}
\]

Step 2: Joint Probabilities

Red number cards:

\[
10 + 10 = 20
\]

\[
Pr(A,B) = \frac{20}{52}
\]

All hearts are red:

\[
Pr(A,C) = \frac{13}{52}
\]

Number cards that are hearts:

\[
Pr(B,C) = \frac{10}{52}
\]

\subsection{Example 2: Costume Party Example}

Alex has:

\begin{itemize}
    \item 3 food-stained white t-shirts
    \item 1 radioactive green cape
    \item 2 dinosaur pajama pants
    \item 4 polka-dot boxer shorts
\end{itemize}

Probability outfit is:

Green cape AND Polka-dot boxers

Options:

\[
\frac{1}{2}, \frac{1}{4}, \frac{1}{6}
\]

Step 1: Analyze Tops

Total tops:

\[
3 + 1 = 4
\]

Probability of cape:

\[
Pr(\text{Cape}) = \frac{1}{4}
\]

(The slide context stops at Step 1, so no further solution steps can be added.)

\section{Conditional Probability}

The lecture outline includes Conditional Probability with Examples 3--5, but the
provided slide context does not include the actual definitions or worked derivations.

Therefore, conditional probability content cannot be written without introducing
material not present in the context.

\section{Exam Revision Summary}

Key formulas from Lecture 4:

\[
Pr(A,B) = Pr(A \cap B)
\]

\[
Pr(A \cup B) = Pr(A) + Pr(B) - Pr(A \cap B)
\]

\[
Pr(A^c) = 1 - Pr(A)
\]

\[
Pr(A,B) = \lim_{n\to\infty} \frac{n_{AB}}{n}
\]

\bigskip

\textbf{End of Lecture 4 Scribe}

\end{document}
